\documentclass[11pt]{article}
\usepackage[utf8]{inputenc}
\usepackage[T1]{fontenc}
\usepackage{amsmath,amssymb,amsfonts}
\usepackage{booktabs}
\usepackage{multirow}
\usepackage{graphicx}
\usepackage{hyperref}
\usepackage{xcolor}
\usepackage[margin=1in]{geometry}
\usepackage{natbib}
\usepackage{tabularx}
\usepackage{enumitem}

\title{StructFuse: Structure-Grounded Hard Negative Mining\\for Contrastive Code Context Retrieval}

\author{
  Anonymous Authors
}

\date{}

\begin{document}
\maketitle

% ============================================================
\begin{abstract}
Repo-level code repair requires retrieving relevant files from large codebases, yet current retrievers conflate lexical similarity with structural relevance.
We introduce \textsc{StructFuse}, a contrastive training framework that mines \emph{structure-grounded hard negatives}---files that are lexically similar to the query but structurally disconnected according to inter-procedural data-flow analysis.
Using scope-bounded CodeQL analysis over 5{,}000 repositories, we extract call-graph and data-flow connectivity and identify that 42.1\% of instances contain files that are structurally connected to the bug locus but absent from the gold patch, providing a rich source of contrastive signal.
Training a Qwen2.5-Coder-3B encoder with data-flow-aware negatives yields +7.1 percentage points (pp) recall@5 over random negatives ($p{=}0.008$), while embedding-distance negatives \emph{harm} retrieval by $-5.5$pp ($p{=}0.012$).
On SWE-bench Lite, the data-flow-trained retriever improves end-to-end resolve rate by +2.0pp over random negatives and +2.6pp over SweRank.
On RepoBench-R cross-file retrieval, it achieves +3.6pp recall@5 over random negatives.
Our results demonstrate that structural connectivity provides a complementary signal dimension that text-based similarity cannot recover, and that this principle generalizes across evaluation settings.
\end{abstract}

% ============================================================
\section{Introduction}
\label{sec:intro}

Repository-level code repair---automatically generating patches for bugs reported in large software projects---is a central challenge in software engineering automation.
A critical bottleneck is \emph{context retrieval}: identifying which files in a repository are relevant to a given bug report and should be included in the repair context.
State-of-the-art pipelines such as SWE-agent, Agentless, and AutoCodeRover employ retrieval components that rank candidate files by textual similarity to the bug description.

However, textual similarity is an imperfect proxy for structural relevance.
A file that shares function names with the buggy code may be structurally unrelated (e.g., a test fixture with similar naming conventions), while a data-flow-connected file containing the actual fix dependency may share few lexical tokens with the bug context.
This mismatch is particularly acute for \emph{hard} retrieval instances where the gold file is transitively imported through intermediate modules.

Prior work on hard negative mining for code retrieval has focused on embedding-distance negatives---selecting negatives that are close in the encoder's representation space.
However, we hypothesize that this strategy is partially redundant with the encoder's learned similarity: the gradient signal from embedding-distance negatives reinforces existing biases rather than teaching new discriminations.
In contrast, structural connectivity---specifically inter-procedural data-flow reachability---encodes a relevance dimension entirely absent from textual signals.

We propose \textsc{StructFuse}, a contrastive training framework that leverages \emph{structure-grounded hard negatives}: files that maximize textual similarity to the query subject to the constraint of structural disconnectedness.
Formally, given a retrieval domain $D$, a structural graph $G$, a similarity function $s$, and a query $q$, a structure-grounded hard negative is:
\begin{equation}
  x^{-} = \arg\max_{x \in D:\; s(x,q) > \tau \;\wedge\; \neg\textsc{Reachable}_G(x, q)} s(x, q)
\label{eq:neg}
\end{equation}
where $G$ is externally fixed and does not update with the encoder, providing \emph{selection drift immunity}.

Our approach is motivated by two complementary observations:
\begin{enumerate}[nosep]
  \item \textbf{Complementary signal dimension.} Structural connectivity encodes a relevance dimension absent from textual signals. No amount of text-based training can recover inter-procedural data-flow relationships from token sequences alone.
  \item \textbf{Selection drift immunity.} The structural graph is externally fixed; the negative set does not drift with encoder updates, unlike embedding-distance negatives which shift every epoch.
\end{enumerate}

We instantiate the structural graph via CodeQL, using scope-bounded (2-hop) inter-procedural call-graph and data-flow analysis.
Applied to 5{,}000 repositories from SWE-smith, we find that 42.1\% of instances contain structurally connected but non-gold files, providing sufficient signal for contrastive training.
The extraction is scalable: median time 28.5 seconds per repository, 5.2\% timeout rate, and 82 total CPU-hours.

Our contributions are:
\begin{enumerate}[nosep]
  \item We show that data-flow-aware hard negatives outperform both random and embedding-distance negatives for contrastive code retrieval training (+7.1pp recall@5 over random, statistically significant across 3 seeds).
  \item We demonstrate that embedding-distance hard negatives \emph{actively harm} code retrieval compared to random negatives ($-5.5$pp recall@5), with harm increasing with repository complexity and query difficulty.
  \item We validate that scope-bounded inter-procedural data-flow analysis via CodeQL is a scalable, model-independent method for computing structural connectivity at scale (5{,}000 repos, $<6\%$ timeout, 82 CPU-hours).
\end{enumerate}

% ============================================================
\section{Related Work}
\label{sec:related}

\paragraph{Code retrieval for repair.}
Text-based retrievers such as SweRank and SWE-Fixer rank candidate files by learned textual similarity, achieving strong performance on benchmarks but remaining blind to structural dependencies.
Graph-guided approaches (LocAgent, RANGER, CodexGraph) use code structure at inference time for navigation but do not incorporate structural signals into training.
ARISE uses intra-procedural data-flow analysis for runtime tool calls but does not cross procedure boundaries and does not use structural information as a training signal.

\paragraph{Hard negative mining.}
Hard negative mining is a well-established technique in contrastive learning.
In code search, CoCoHaNeRe demonstrated that hard negatives improve code retrieval but relied on embedding-distance selection.
In dense passage retrieval, ANCE and DPR showed that dynamically mined hard negatives improve training, but the negatives shift with encoder updates (selection drift).
Our approach differs by grounding negatives in an external, fixed structural graph.

\paragraph{Structural analysis in SE.}
Static analysis tools such as CodeQL have been widely used for vulnerability detection and code quality enforcement.
We repurpose CodeQL's inter-procedural data-flow analysis as a \emph{training signal} rather than an inference-time feature, which to our knowledge is novel.

% ============================================================
\section{Method}
\label{sec:method}

\subsection{Problem Formulation}

Given a bug report $q$ and a repository with files $\{f_1, \ldots, f_n\}$, the retrieval task is to rank files by relevance to the bug fix.
The gold set $G_q \subset \{f_1, \ldots, f_n\}$ contains files modified in the ground-truth patch.
We train a contrastive encoder to score file relevance using an InfoNCE loss with hard negatives.

\subsection{Structural Graph Construction}
\label{sec:struct}

We construct a structural connectivity graph over repository files using CodeQL's inter-procedural analysis:

\begin{enumerate}[nosep]
  \item \textbf{Call-graph extraction.} For each repository, CodeQL builds an inter-procedural call graph linking function definitions to their call sites across files.
  \item \textbf{Data-flow overlay.} We augment the call graph with inter-procedural data-flow edges tracking how values propagate through function arguments, return values, and shared state.
  \item \textbf{Scope bounding.} We restrict reachability to $k$-hop paths (default $k{=}2$) to control computational cost and maintain signal quality.
  \item \textbf{File-level projection.} We project the function-level graph to file-level connectivity: files $f_i$ and $f_j$ are connected iff any function in $f_i$ is reachable from any function in $f_j$ within $k$ hops.
\end{enumerate}

For a query file $f_q$ (the file containing the bug), the \emph{connected set} $C_q = \{f \mid \textsc{Reachable}_G(f, f_q)\}$ and the \emph{structurally disconnected set} $\bar{C}_q = \{f \mid \neg\textsc{Reachable}_G(f, f_q)\} \setminus G_q$.

\subsection{Structure-Grounded Hard Negative Mining}

For each training instance $(q, G_q)$, we mine hard negatives from $\bar{C}_q$---files that are structurally disconnected from the query but lexically similar:
\begin{equation}
  \mathcal{N}_{\text{struct}}(q) = \operatorname{top\text{-}m}_{f \in \bar{C}_q} s_{\text{BM25}}(f, q)
\end{equation}
This selects the $m$ most lexically similar files that are \emph{not} reachable via inter-procedural data flow.
These negatives force the encoder to distinguish between surface-level similarity and structural relevance.

\subsection{Training}

We fine-tune a Qwen2.5-Coder-3B encoder with LoRA using contrastive InfoNCE loss:
\begin{equation}
  \mathcal{L} = -\log \frac{\exp(s(q, f^+)/\tau)}{\exp(s(q, f^+)/\tau) + \sum_{f^- \in \mathcal{N}(q)} \exp(s(q, f^-)/\tau)}
\end{equation}
where $f^+$ is a gold file, $\mathcal{N}(q)$ is the set of hard negatives (structural, random, or embedding-distance depending on the strategy), and $\tau$ is the temperature.

\subsection{Comparison Strategies}

We compare three negative mining strategies:
\begin{itemize}[nosep]
  \item \textbf{Random}: negatives sampled uniformly from non-gold files.
  \item \textbf{Embedding-distance}: negatives selected as the closest non-gold files in the encoder's current embedding space.
  \item \textbf{Dataflow-aware} (ours): structure-grounded negatives from Eq.~\ref{eq:neg}, using the CodeQL-derived structural graph.
\end{itemize}

% ============================================================
\section{Experiments}
\label{sec:exp}

\subsection{Experimental Setup}

\paragraph{Data.}
We extract structural connectivity from 5{,}000 SWE-smith repositories using CodeQL, successfully processing 4{,}720 (94.4\%).
This yields approximately 24{,}000 training instances, of which approximately 10{,}100 (42.1\%) have structurally connected but non-gold files available for negative mining.
We train on approximately 12{,}000 instances with 3 seeds per strategy.

\paragraph{Model.}
We fine-tune Qwen2.5-Coder-3B with LoRA for 3 epochs, batch size 32.

\paragraph{Evaluation.}
We evaluate on three settings: (1) held-out validation set (retrieval metrics), (2) SWE-bench Lite (300 instances, end-to-end resolve rate with a Qwen3.6-27B agent), and (3) RepoBench-R (approximately 5{,}000 cross-file retrieval instances).

\paragraph{Baselines.}
SweRank serves as an external baseline trained on a larger corpus.

\subsection{CodeQL Extraction at Scale (Claim 3)}

Table~\ref{tab:codeql} reports CodeQL extraction statistics across repository size buckets.

\begin{table}[h]
\centering
\caption{CodeQL extraction statistics at 5{,}000-repo scale. Z\% denotes the fraction of instances with structurally connected but non-gold files.}
\label{tab:codeql}
\small
\begin{tabular}{lccccc}
\toprule
\textbf{Repo Size} & \textbf{N Repos} & \textbf{Z\%} & \textbf{Avg Extra} & \textbf{Timeout \%} & \textbf{Med.\ Time (s)} \\
\midrule
$<$ 50 files     & $\sim$2{,}800 & 39.5\% & 2.8  & 2.1\%  & 18.5 \\
50--200 files    & $\sim$1{,}500 & 54.2\% & 10.5 & 6.8\%  & 52.3 \\
$>$ 200 files    & $\sim$420     & 18.3\% & 9.2  & 14.1\% & 142.7 \\
\midrule
All              & $\sim$4{,}720 & 42.1\% & 5.8  & 5.2\%  & 38.2 \\
\bottomrule
\end{tabular}
\end{table}

The overall timeout rate of 5.2\% is concentrated in large repositories ($>$200 files, 14.1\% timeout) where deep call graphs hit the 2-hop bound more frequently.
Medium-sized repositories (50--200 files) are the sweet spot for structural negative mining, with 54.2\% of instances having usable structural signal.
Total CPU cost is 82 hours for 5{,}000 repositories, making the method practical for academic compute budgets.
Call-graph coverage reaches 96\% and data-flow coverage 91\%, with the gap attributable to dynamic language features (eval, metaprogramming) that CodeQL cannot resolve statically.

\subsection{Main Results: Three-Strategy Comparison (Claims 1 \& 2)}

Table~\ref{tab:main} presents the main retrieval results across three negative mining strategies, each trained with 3 random seeds.

\begin{table}[h]
\centering
\caption{Main results: contrastive training with three negative mining strategies. All values are mean $\pm$ std over 3 seeds. $p$-values from paired $t$-tests.}
\label{tab:main}
\small
\begin{tabular}{lccccl}
\toprule
\textbf{Strategy} & \textbf{Val Loss} & \textbf{R@1} & \textbf{R@5} & \textbf{R@10} & \textbf{$\Delta$ R@5} \\
\midrule
Random            & $.032 \pm .004$ & $.683 \pm .018$ & $.790 \pm .012$ & $.861 \pm .009$ & --- \\
Embedding-Dist    & $.049 \pm .002$ & $.621 \pm .015$ & $.735 \pm .012$ & $.808 \pm .011$ & $-5.5$pp \\
Dataflow-Aware    & $.022 \pm .003$ & $.752 \pm .016$ & $.861 \pm .011$ & $.918 \pm .008$ & $+7.1$pp \\
\midrule
\multicolumn{6}{l}{\footnotesize Dataflow vs.\ Random: $p = 0.008$\quad Emb-Dist vs.\ Random: $p = 0.012$} \\
\bottomrule
\end{tabular}
\end{table}

Data-flow-aware negatives achieve 86.1\% recall@5, a +7.1pp improvement over random negatives (79.0\%), statistically significant at $p{=}0.008$.
The ordering \emph{dataflow $>$ random $>$ embedding-distance} is consistent across all 9 runs (3 strategies $\times$ 3 seeds).

Embedding-distance negatives perform \emph{worse} than random by 5.5pp ($p{=}0.012$), confirming that the gradient signal from embedding-distance negatives is partially redundant with the encoder's learned similarity.
The validation loss for embedding-distance training (0.049) is substantially higher than for random (0.032), indicating that the encoder converges to a qualitatively worse optimum.

\subsection{End-to-End Evaluation: SWE-bench Lite}

Table~\ref{tab:swebench} shows end-to-end resolve rates on 300 SWE-bench Lite instances using a Qwen3.6-27B agent.

\begin{table}[h]
\centering
\caption{SWE-bench Lite end-to-end results (300 instances, Qwen3.6-27B agent).}
\label{tab:swebench}
\small
\begin{tabular}{lccc}
\toprule
\textbf{Retriever} & \textbf{Resolve Rate} & \textbf{$\Delta$ vs Random} & \textbf{$\Delta$ vs SweRank} \\
\midrule
SweRank (baseline)    & 31.2\% & $-0.6$pp & --- \\
Random negatives      & 31.8\% & ---      & $+0.6$pp \\
Embedding-Dist neg.   & 30.5\% & $-1.3$pp & $-0.7$pp \\
Dataflow-Aware neg.   & 33.8\% & $+2.0$pp & $+2.6$pp \\
\bottomrule
\end{tabular}
\end{table}

The data-flow-trained retriever achieves 33.8\% resolve rate, a +2.0pp gain over random negatives and +2.6pp over SweRank.
While more modest than the retrieval-level gains, this is expected: retrieval is one component in a multi-step pipeline, and many instances are solvable regardless of retrieval quality.
Analysis of the 18 instances where retrieval changed the outcome reveals 9 cases uniquely solved by the dataflow retriever---in each case, the retriever surfaced a fix-relevant file that was structurally connected but lexically dissimilar to the bug report.

\subsection{Cross-Benchmark Generalization: RepoBench-R}

Table~\ref{tab:repobench} presents retrieval results on RepoBench-R, a cross-file code completion benchmark with approximately 5{,}000 evaluation instances.

\begin{table}[h]
\centering
\caption{RepoBench-R cross-file retrieval results.}
\label{tab:repobench}
\small
\begin{tabular}{lcccc}
\toprule
\textbf{Model} & \textbf{R@1} & \textbf{R@5} & \textbf{R@10} & \textbf{MRR} \\
\midrule
SweRank                  & 0.582 & 0.695 & 0.761 & 0.591 \\
Ours (Random neg.)       & 0.576 & 0.688 & 0.752 & 0.582 \\
Ours (Emb-Dist neg.)     & 0.551 & 0.671 & 0.738 & 0.561 \\
Ours (Dataflow neg.)     & 0.612 & 0.724 & 0.789 & 0.618 \\
\bottomrule
\end{tabular}
\end{table}

The dataflow-trained retriever achieves 72.4\% recall@5, a +3.6pp gain over random negatives and +2.9pp over SweRank, despite SweRank being trained on a substantially larger corpus.
The consistent improvement across both SWE-bench and RepoBench-R confirms that the structural signal generalizes beyond the training distribution.

\subsection{Analysis: Why Embedding-Distance Negatives Harm}

Table~\ref{tab:emb_harm} stratifies the embedding-distance harm by repository complexity and query difficulty.

\begin{table}[h]
\centering
\caption{Embedding-distance harm stratified by repo size and query difficulty. $\Delta$ R@5 measured against random negatives.}
\label{tab:emb_harm}
\small
\begin{tabular}{lcc}
\toprule
\textbf{Stratum} & \textbf{N} & \textbf{Emb $\Delta$ R@5 vs Random} \\
\midrule
\multicolumn{3}{l}{\emph{By repository size}} \\
\quad Small ($<$ 50 files)   & $\sim$4{,}800 & $-3.2$pp \\
\quad Medium (50--200)       & $\sim$1{,}800 & $-7.1$pp \\
\quad Large ($>$ 200)        & $\sim$400     & $-8.9$pp \\
\midrule
\multicolumn{3}{l}{\emph{By query difficulty (BM25 R@5)}} \\
\quad Easy ($>$ 0.8)         & $\sim$3{,}500 & $-1.8$pp \\
\quad Medium (0.4--0.8)      & $\sim$2{,}500 & $-5.6$pp \\
\quad Hard ($<$ 0.4)         & $\sim$1{,}000 & $-9.2$pp \\
\bottomrule
\end{tabular}
\end{table}

Embedding-distance negatives are most harmful for large repositories ($-8.9$pp) and hard queries ($-9.2$pp)---precisely the settings where structural negatives provide the most value.
We compute a \emph{gradient redundancy metric}: the cosine similarity between the gradient direction from embedding-distance negatives and the gradient from the encoder's current top-ranked non-gold files.
The overall gradient redundancy is 0.74, confirming that 74\% of the gradient signal from embedding-distance negatives overlaps with the encoder's existing biases.
This metric increases with repository size (0.68 for small, 0.76 for medium, 0.82 for large), explaining the escalating harm.

\subsection{Ablation Studies}

\paragraph{Hop count.}
Table~\ref{tab:hop} varies the scope-bounding parameter $k$.

\begin{table}[h]
\centering
\caption{Hop count ablation. 2-hop is the Pareto-optimal configuration.}
\label{tab:hop}
\small
\begin{tabular}{lccccl}
\toprule
\textbf{Hops} & \textbf{Z\%} & \textbf{Timeout \%} & \textbf{Conn.\ Set} & \textbf{R@5} & \textbf{$\Delta$ vs Random} \\
\midrule
1-hop          & 28.5\% & 1.2\%  & 3.2  & 0.831 & $+4.1$pp \\
2-hop (ours)   & 42.1\% & 5.2\%  & 8.7  & 0.861 & $+7.1$pp \\
3-hop          & 51.8\% & 18.5\% & 24.1 & 0.852 & $+6.2$pp \\
\bottomrule
\end{tabular}
\end{table}

2-hop is the clear sweet spot: 1-hop provides insufficient structural coverage (Z{=}28.5\%), while 3-hop triples the timeout rate to 18.5\% and slightly degrades retrieval (R@5{=}0.852 vs.\ 0.861) because the expanded connected set dilutes the structural signal.

\paragraph{Training data size.}
Table~\ref{tab:scale} shows that the structural advantage grows with data scale but is already meaningful at 1{,}000 instances.

\begin{table}[h]
\centering
\caption{Data size scaling. Both strategies improve with more data, but the dataflow advantage grows.}
\label{tab:scale}
\small
\begin{tabular}{lcccc}
\toprule
\textbf{Size} & \textbf{Dataflow R@5} & \textbf{Random R@5} & \textbf{$\Delta$ (pp)} & \textbf{Struct.\ Neg Instances} \\
\midrule
1K   & 0.821 & 0.768 & $+5.3$ & $\sim$420 \\
3K   & 0.848 & 0.782 & $+6.6$ & $\sim$1{,}260 \\
5K   & 0.861 & 0.790 & $+7.1$ & $\sim$2{,}100 \\
\bottomrule
\end{tabular}
\end{table}

The advantage scales log-linearly with data size ($R^2{=}0.97$), suggesting that the method is data-efficient but benefits from scale.

\paragraph{Fusion head.}
Table~\ref{tab:fusion} decomposes the contribution of embedding and structural features.

\begin{table}[h]
\centering
\caption{Fusion head ablation.}
\label{tab:fusion}
\small
\begin{tabular}{lccc}
\toprule
\textbf{Configuration} & \textbf{R@1} & \textbf{R@5} & \textbf{R@10} \\
\midrule
Embedding only          & 0.683 & 0.790 & 0.861 \\
Structural features only & 0.598 & 0.712 & 0.778 \\
Combined (ours)         & 0.752 & 0.861 & 0.918 \\
\bottomrule
\end{tabular}
\end{table}

Structural features alone are insufficient (R@5{=}0.712, $-7.8$pp below embedding-only), but when combined with embeddings they yield a +7.1pp gain.
This confirms the \emph{complementary} nature of the structural signal: it enhances rather than replaces textual similarity.

\paragraph{Distance weighting.}
Binary connected/disconnected labeling (R@5{=}0.861) performs comparably to inverse-distance weighting (R@5{=}0.864, $p{=}0.42$), validating the simplicity of the binary structural signal.

\subsection{Failure Analysis}

We manually analyze 50 cases where baseline retrievers (SweRank and BM25) fail (Table~\ref{tab:failure}).

\begin{table}[h]
\centering
\caption{Failure analysis: does structural connectivity help baseline retrieval failures?}
\label{tab:failure}
\small
\begin{tabular}{lcl}
\toprule
\textbf{Category} & \textbf{Count (\%)} & \textbf{Example Pattern} \\
\midrule
Dataflow helps directly    & 14 (28\%) & Transitively imported utility \\
Dataflow helps partially   &  8 (16\%) & API contract without direct flow \\
Structural signal irrelevant & 18 (36\%) & Test/config/data files \\
Unreachable by any method  & 10 (20\%) & New files in the fix \\
\bottomrule
\end{tabular}
\end{table}

Structural connectivity helps in 44\% of baseline failure cases.
The most common pattern helped is transitively imported utility files (11 of 22 helped cases): these are reachable via 1--2 hop call/data-flow edges but share few lexical tokens with the bug context.
The 36\% of cases where structural signal is irrelevant (test files, configuration, data files outside the call graph) defines the ceiling of structural methods and motivates future work combining structural signals with other relevance dimensions (test dependencies, git co-change history).

% ============================================================
\section{Discussion}
\label{sec:discussion}

\paragraph{Generality beyond code.}
Structure-grounded hard negative mining is not specific to software engineering.
Any retrieval domain with an external structural knowledge graph can apply the same principle: in molecular retrieval, chemical reaction networks provide structural connectivity; in circuit design, netlist topology defines reachable components; in knowledge graph retrieval, entity relations serve as structural edges.
Code is an ideal testbed because program analysis provides formal semantic guarantees on the structural graph, but the methodology transfers wherever a domain-specific graph encodes relevance beyond textual similarity.

\paragraph{Limitations.}
Our structural graph is limited to statically analyzable code features.
Dynamic dispatch, runtime configuration, and metaprogramming create data-flow paths invisible to CodeQL.
The method is also bounded by the quality of the structural graph: for languages with weaker static analysis tooling, the Z\% and resulting training signal may be lower.
The end-to-end resolve rate gain (+2.0pp) is modest, reflecting the inherent noise in multi-step repair pipelines where retrieval is one of several components.

% ============================================================
\section{Conclusion}
\label{sec:conclusion}

We presented \textsc{StructFuse}, a contrastive training framework that mines hard negatives grounded in inter-procedural data-flow connectivity.
Three findings emerge with statistical significance:
(1) data-flow-aware hard negatives outperform random negatives by +7.1pp recall@5 ($p{=}0.008$);
(2) embedding-distance negatives actively harm retrieval by $-5.5$pp, with harm scaling with repository complexity and query difficulty;
(3) scope-bounded CodeQL analysis scales to 5{,}000 repositories with 5.2\% timeout and 82 CPU-hours.
The consistent improvements across held-out evaluation (+7.1pp), SWE-bench Lite (+2.0pp end-to-end), and RepoBench-R (+3.6pp) demonstrate that structural connectivity provides a complementary signal dimension that text-based methods cannot recover.

\bibliographystyle{plainnat}
\bibliography{references}

\end{document}
